\chapter{Methods} % Chapter 3
\label{Methods}

% =============================================================================
% METHODS (10-15 pages)
% Rubric: Methods (20% of Report)
% "Were the methods described in such a way that the study could be repeated
%  and that the results could be understood properly? Was a justification
%  given for the choice of methods?"
% =============================================================================

% -----------------------------------------------------------------------------
% SECTION 1: Experimental Design Overview
% -----------------------------------------------------------------------------
\section{Experimental Design}
% TODO: High-level description of comparative methodology
% TODO: Justify comparative approach - why compare RL vs LLM?
% TODO: Describe the three conditions (from proposal Table 1):
%       1. RL (LSTM) - baseline
%       2. LLM-Comparison - same information as RL
%       3. LLM-Control - prompt variation for robustness

\subsection{Controlled Comparison Rationale}
% TODO: What variables are held constant across conditions?
%       - Environment (game rules)
%       - Available information
%       - Optimization objective (resource maximization)
% TODO: What differs? Only the agent architecture
% TODO: Acknowledge: cannot isolate ALL architectural differences

% -----------------------------------------------------------------------------
% SECTION 2: Game Design
% -----------------------------------------------------------------------------
\section{Coordination Game Design}
% TODO: Justify game choice - why this minimal coordination game?
% TODO: Game mechanics from proposal:
%       - Agents pursue resource optimization
%       - Simultaneous action selection each round

\subsection{Action Space}
% TODO: Define all actions (from proposal Table 2):
%       - Invest (self/other): costs V, generates I
%       - Arm (self/other): costs D, multiplies M
%       - Attack (other): resource relocation
% TODO: Why these specific actions? What social structures can they produce?

\subsection{Conflict Resolution}
% TODO: Probabilistic resolution based on combined resources
% TODO: Winner takes percentage T of loser's resources
% TODO: All combatants pay conflict costs C
% TODO: Specify all parameter values (V, I, D, M, T, C, N)

\subsection{Information Structure}
% TODO: Complete information - all resources visible
% TODO: History window for tracking previous interactions
% TODO: N rounds (unknown to players) - why?

% -----------------------------------------------------------------------------
% SECTION 3: Agent Implementations
% -----------------------------------------------------------------------------
\section{Agent Architectures}

\subsection{RL Agent (LSTM)}
% TODO: Policy network architecture
% TODO: LSTM for temporal information (addressing Markovian limitation)
% TODO: Training procedure - how many episodes, reward function
% TODO: Hyperparameters

\subsection{LLM Agent}
% TODO: Which LLM (model, version, API)?
% TODO: Prompt structure for comparison condition
%       - What information is provided?
%       - Goal statement (optimize resources)
% TODO: Prompt variations for control condition
%       - Tone, decoration, language, explicit instructions

% -----------------------------------------------------------------------------
% SECTION 4: Metrics
% -----------------------------------------------------------------------------
\section{Measurement of Emergent Social Structure}
% TODO: Operationalize "emergent social structure"
% TODO: Justify each metric choice

\subsection{Coalition Structure}
% TODO: Community detection algorithm (which one?)
% TODO: Directed interaction networks from investment/arm actions
% TODO: Metrics: number of communities, average size, temporal stability

\subsection{Interaction Patterns}
% TODO: Cooperation frequency (investment actions)
% TODO: Aggression frequency (attack actions)
% TODO: Retaliation probability - how computed?

\subsection{Hierarchy}
% TODO: Gini coefficient over resource distribution
% TODO: Temporal dynamics of inequality

% -----------------------------------------------------------------------------
% SECTION 5: Experimental Protocol
% -----------------------------------------------------------------------------
\section{Experimental Protocol}
% TODO: How many simulation runs per condition?
% TODO: Number of agents per simulation
% TODO: Number of rounds per simulation
% TODO: Statistical tests for comparing conditions
% TODO: Sensitivity analysis plan

% -----------------------------------------------------------------------------
% SECTION 6: Implementation Details
% -----------------------------------------------------------------------------
\section{Implementation}
% TODO: Programming language, libraries, frameworks
% TODO: Computational resources used
% TODO: Code availability statement
% NOTE: Rubric asks "Did the student demonstrate programming skills?"
